\section{Discussion and Limitations}\label{sec:discussion}

\subsection{Label Resolution, Not Sensing}
Under the fixed footprint cut, 90 drive-level macro-F1 needs each held-out model's footprint to within about $\pm$0.12~m$^2$ (Fig.~\ref{fig:ceiling}). Two curb-mass cuts fitted on training models score only 63.3, and 21 of the 41 held-out car models lie within 0.15~m$^2$ of a differently labelled training model. The gap is set by label resolution and fleet composition, not by model capacity or evidence per drive: architectures stay within about 2 points (Table~\ref{tab:main}) and more windows per drive add little for size (Fig.~\ref{fig:evidence}). Per-model errors (Section~\ref{sec:experiments}) split LARGE by body type: pickups and body-on-frame SUVs are recognised (GMC Yukon, Ford Ranger and Ford F-150 Lightning: 95.5--100\% of drives), while unibody three-row SUVs and full-size sedans are read as MID or SMALL (Genesis G80 and G90, Volkswagen Atlas, Tesla Model~S and Kia EV9: 0--5.6\%), so the networks appear to learn LARGE as a truck's mass-and-driveline signature rather than as footprint.

\subsection{What 90 Needs}
Pooling a vehicle's drives does not rescue size (60.0--62.6 against 66.8--67.1 at drive level): about 45\% of the test vehicles have a single drive and the errors are per car model, so even the per-car-model majority oracle lifts the reference ensemble only from 63.2 to 72.0. Three routes remain: seen car models, as under held-out devices, where the same recipes reach 79.5--80.6 (Table~\ref{tab:device}); the vehicle level for BEV vs ICE, where the GRU B5 rises from 92.6 to 93.4 (95.5 on vehicles with at least two drives) and the ResNet B5 stays at 92.8--93.0; and labels whose margins exceed the resolution of about 0.12~m$^2$, such as cuts in wide gaps of the footprint distribution.

\subsection{Two Protocol Lessons}
First, while a driving automation system (DAS) is engaged, a controller tuned per car model supplies part of the control action. Windows with automation activity are easier for size than windows without (56.4 against 50.2, one bi-GRU run) and harder for BEV vs ICE (79.8 against 84.0, one ResNet B5 run): plausibly the controller carries platform information and masks the pedal behaviour that separates the powertrains. Hence the headline drive-level results (Table~\ref{tab:main}) are paired with the human-only protocol, and the final BEV vs ICE recipe (B5) down-weights DAS windows (the size variant S7 did not win in the fold-0 sweep, Table~\ref{tab:folds}). Second, early stopping on the released validation split, which shares car models with training, selects epochs 16--22, against 1--8 (median about 4) for the reference recipe on the model-disjoint folds (Section~\ref{sec:methods}), and the final models' validation scores do not predict their test scores (Section~\ref{sec:experiments}). Such a split rewards memorising car models, a shortcut~\cite{geirhos2020shortcut}; we therefore fixed the training length from the folds (Table~\ref{tab:folds}), and a future release would carve the validation split model-disjoint.

\subsection{Privacy and Limitations}
Section~\ref{sec:dataset} lists the privacy measures. Labels describe car models, not people: a rare model may be recognisable as a model, never as a person. The same property that makes the data useful, that motion reveals the vehicle, restricts access to research on vehicle dynamics and driving behaviour and prohibits re-identification and linkage with other datasets. Several limitations remain.
All logs come from one logger ecosystem and its software eras~\cite{schafer2018comma2k19}, and the fleet is skewed by region and brand (83.8\% of the hours from North America, 4.8\% from Europe, 4.2\% from Asia); region, software era and device type are metadata, not inputs, but they correlate with powertrain.
Labels are platform attributes: a misidentified platform mislabels every drive of that vehicle, and platforms sold with several powertrains the logs cannot separate, or of unknown powertrain, carry no powertrain label (11.1\% of the drives). On 56.2\% of the drives the lead vehicle comes from a camera model rather than a radar, and pedal magnitudes are missing on 63.9\% (accelerator) and 49.4\% (brake) of the drives. The internal folds never hold out the LARGE~$\times$~BEV cell, which one car model dominates, so their numbers for that cell reflect training models only; the test split does hold out models of that cell.
% Source of the fold statement: CAMPAIGN.md (G1/G2 consequence note) and INTERNAL_FOLDS.md (the LARGE x BEV pool is dominated by one car model); not in OUTLINE.md; verified 8 Nov 2026.
With 41 (size) and 39 (BEV vs ICE) test car models, the GRU recipes' cluster-bootstrap intervals are wide (67.1 [54.4, 74.2] and 92.6 [87.1, 96.2]); differences of a few points between methods lie inside them.
